\documentclass[12pt]{article}
\usepackage[utf8]{inputenc}
\usepackage[T1]{fontenc}
\usepackage{hyperref}
\title{Nahuat Language Resources for K-12 Curriculum Development}
\author{Compiled by Claude Sonnet 5.5 (claude-sonnet-5-5) for Andr\'as Kornai}
\date{}
\begin{document}
\maketitle

\section{Introduction}
Nahuat (Nawat, Pipil; ISO 639-3 ppl) is a Uto-Aztecan language of El Salvador. Estimates of speakers vary widely, from fewer than 100 fluent speakers (Campbell, King and Lemus 2014) to figures of several hundred; the reference figure of about 10,000 is likely an overestimate or counts ethnic population.

\section{Primary Educational Resources}
Search summaries report a Don Bosco University revitalization programme teaching Nawat in public schools (started 2003 with 3 schools; later reported as 11 schools), and a series of texts produced for schools in Izalco and Nahuizalco. The primary pages were not retrieved, so details are as reported in search summaries.
\begin{itemize}
\item Endangered Languages Project, Nawat: \url{https://endangeredlanguages.com/lang/481?hl=en}
\item Global Voices Rising, ``A Video Library for Successor Pipil Generation'': \url{https://rising.globalvoices.org/blog/2012/08/27/languages-a-video-library-for-successor-pipil-generation/}
\end{itemize}

\section{Linguistic and Reference Materials}
\begin{itemize}
\item Wikipedia, Nawat language: \url{https://Www.wikipedia.org/wiki/Nawat}
\item Article in the journal \emph{Di\'alogos} (CAMJOL), surfaced in the search on Nawat teaching; its content was not examined: \url{https://camjol.info/index.php/DIALOGOS/article/view/16312/19289}
\end{itemize}
Alan King's elementary Nawat course (2011) and a Living Tongues talking dictionary were mentioned in search summaries but no direct URLs were seen.

\section{Religious and Community Resources}
None found in search results.

\section{Digital and Audio Resources}
The Global Voices item above concerns a video library for Pipil learners.

\section{Curriculum Development Context}
Nawat has an existing school-based revitalization effort with non-native teachers trained at the university, which is the most relevant precedent for curriculum work.
\end{document}